\label{sec:bootstrap}
Khalkhali and Pagliaroli \cite{KP} solve Dirac ensembles without solving them: Schwinger--Dyson (loop) equations give
linear relations among moments, positivity of moment matrices cuts out a convex island, and the island shrinks onto
the answer as the moment degree grows. The network admits the same treatment, with the Gaussian input as the
``action'', Stein's identity as the loop equation, and the gates as idempotent variables.

\begin{definition}[The ReLU moment problem]\label{def:moment}
Let $\R[x,\theta]$ be the polynomial algebra in $x\in\R^n$ and commuting idempotents $\theta_{l,a}$ ($\theta^2=\theta$),
$l=1,\dots,L$. Define recursively the polynomials $z_1=W_1x$, $z_{l,a}=\sum_bW_l(a,b)\,\theta_{l-1,b}\,z_{l-1,b}$. A
\emph{feasible moment functional} is a linear $\mathcal L:\R[x,\theta]\to\R$ with $\mathcal L(1)=1$, Gaussian
$x$-moments (equivalently the loop equations $\mathcal L(x_iq)=\mathcal L(\partial_iq)$ for $q\in\R[x]$), positive moment
matrix $\mathcal L(p^2)\ge0$, and positive localizing forms $\mathcal L(p^2\,\theta_{l,a}z_{l,a})\ge0$,
$\mathcal L(-p^2(1-\theta_{l,a})z_{l,a})\ge0$ for all $p$.
\end{definition}

\begin{theorem}[The island is a point]\label{thm:island}
For almost every weight realisation every feasible moment functional represented by a measure is the law of
$(x,\theta(x))$ except for the values of the gates on dead tiles, where all downstream pre-activations vanish and the
gates are unconstrained; these measures are moment-determinate, and they all give the same value to every
pre-activation moment. In particular $u=\mathcal L(\theta_L\circ z_L)$ is determined by positivity and the loop
equations alone.
\end{theorem}
\begin{proof}
A representing measure $\mu$ lives on $\R^n\times\{0,1\}^{nL}$ and is supported in
$K=\{\theta_{l,a}z_{l,a}\ge0,\ (1-\theta_{l,a})z_{l,a}\le0\}$. If $z_{1,a}(x)>0$ the second constraint forces
$\theta_{1,a}=1$, if $z_{1,a}(x)<0$ the first forces $\theta_{1,a}=0$; given the gates of layers $<l$, $z_l(x,\theta)$ is the
true pre-activation and the same argument forces layer $l$. So the fibre of $K$ over $x$ is the single point $\theta(x)$
unless some $z_{l,a}(x)=0$. On a live tile this happens on a $\gamma$-null set ($z_{l,a}$ is a nonzero linear form there
almost surely); on a dead tile, where $h_l\equiv0$, all later pre-activations vanish whatever the later gates are, so
the fibre is a set of gate patterns on which every $z_{l'}$, $l'>l$, and $h_L$ take the value $0$. Since the $x$-marginal
is $\gamma$, $\mu$ is the graph measure up to the choice of these irrelevant gates. Determinacy: $\mu=\sum_\vartheta\mu_\vartheta\otimes\delta_\vartheta$
over the finitely many gate patterns, with $\mu_\vartheta\le\gamma$; the moments of $\mu$ determine those of every $\mu_\vartheta$
(the monomials in $\theta$ span the indicators of patterns), and each $\mu_\vartheta$ satisfies Carleman's condition because its
moments are bounded by the Gaussian ones.
\end{proof}

Under the standard convergence theory of moment--sum-of-squares relaxations for determinate measures
\cite{Las}, the islands of degree $2k$ form a nested sequence of intervals containing $u$ and shrinking to it. Two
features are specific to the network.

\begin{proposition}[Loop equations need the kink generators]\label{prop:loops}
Stein's identity for test functions that involve gates produces the distributions $\delta(z_{l,a})$:
$\partial_i\theta_{l,a}=\delta(z_{l,a})\,\partial_iz_{l,a}$. The loop equations therefore close only on the algebra generated by
$x$, the gates and the kink distributions $\delta^{(k)}(z_{l,a})$. The first of them, for the test field $q=J_{\cdot i}$, is
the kink-current formula \eqref{eq:kink}; the Wick coefficients $\E[\partial^k\relu^p]$ of every cumulant chain are the
Gaussian-state values of these kink generators.
\end{proposition}
\begin{proof}
Differentiate $\theta_{l,a}(x)=\1[z_{l,a}(x)>0]$ in the sense of distributions; the first loop equation is
$\E[x\cdot\nabla h_L]=\E[\Delta h_L]$, computed in Theorem~\ref{thm:kink}. For a Gaussian $z$, $\E[\delta^{(k)}(z)]$ are the
derivatives of the density at the kink, which is what the chains' Wick tables contain.
\end{proof}

\begin{proposition}[The level-one island of a single neuron is Scarf's interval]\label{prop:scarf}
If only the mean $m$ and variance $s^2$ of $z_{l,a}$ are imposed, together with $h\ge0$, $h\ge z$, $h(h-z)=0$, the closure of the set of feasible
values of $\E h$ is exactly $\big[\max(0,m),\ \tfrac12(m+\sqrt{m^2+s^2})\big]$. The Gaussian closure picks the interior
point $s\,(\alpha\Phi(\alpha)+\varphi(\alpha))$, $\alpha=m/s$.
\end{proposition}
\begin{proof}
The constraints force $h=z_+$; the lower bound is Jensen, and the upper bound is Scarf's worst case over laws with
given mean and variance, attained by a two-point law \cite{Scarf}.
\end{proof}

So a moment chain is a truncated Schwinger--Dyson hierarchy solved with a Gaussian (or Edgeworth) ansatz and
\emph{without} positivity. Positivity is what would turn a chain into a certificate; its cost at width $1024$ is
prohibitive at any useful degree, and its role here is structural: it shows that the exact answer is the unique point
of an island cut out by the input's loop equations and the complementarity of the gates, and that the information a
chain discards is the positivity of the kink generators' moment matrices, which is exactly what goes wrong when a
closure produces a negative fourth cumulant or a variance outside its admissible range.
